\section*{Methods}
\subsection*{Channel, feasible decisions and opportunity loss}
We consider a base station with \(N\) jointly processed antenna streams serving \(K=2S\) users on \(S\) labelled, orthogonal tones of equal bandwidth. Each user transmits one independent unit-variance proper complex Gaussian symbol on its assigned tone. For the two users \(i,j\) on tone \(s\), the whitened observation is \(y_s=\sqrt{p_i}h_{is}x_i+\sqrt{p_j}h_{js}x_j+n_s\), where the rate model takes \(n_s\) to be proper complex Gaussian with covariance \(I_N\), independent of the users' symbols. In simulation, powers satisfy \(0.05\leq p_i\leq0.20\) W and \(\sum_i p_i\leq P=0.10K\) W. The Rice recordings use digital receiver units. The oracle serves every user once, optimizes within these power limits and uses a separate linear combiner for each user, treating its partner as interference.

For nonzero \(h_{is}\), Cauchy--Schwarz in the interference-plus-noise inner product gives the optimal combiner direction \((I_N+p_j h_{js}h_{js}^{H})^{-1}h_{is}\). Let \(x=p_i\|h_{is}\|^2\) and \(y=p_j\|h_{js}\|^2\). The squared normalized overlap \(r=|h_{is}^{H}h_{js}|^2/(\|h_{is}\|^2\|h_{js}\|^2)\) is set to zero if either channel vanishes. The received signal-to-interference-plus-noise ratios and pair rate in nats are
\begin{equation}
\begin{aligned}
\gamma_i &= x\left(1-\frac{r y}{1+y}\right), &
\gamma_j &= y\left(1-\frac{r x}{1+x}\right),\\
q_{e,s} &= \ln(1+\gamma_i)+\ln(1+\gamma_j), & e&=\{i,j\}.
\end{aligned}
\label{eq:sinr}
\end{equation}
The oracle rate \(V(h)\) maximizes the sum of pair rates over a finite set of schedules and a compact set of powers. Evaluating the stored filters at the implemented powers gives the Gaussian-input mutual information \(\fix(h)\), assuming that the decoder knows the effective scalar complex gain \(v_i^Hh_{is}\) after combining, even when those filters are suboptimal. On any channel for which both endpoints hold, \(U-L=(U-V)+(V-\fix)+(\fix-L)\). Only the middle term is the rate a better feasible decision would recover. If several methods bound the same \(V\), their smallest checked ceiling also bounds it. We paired this ceiling \(U_{\min}\) with each method's floor on the same known channel. For \(V>0\) and \(L\ge0\), \(L/U_{\min}\le\fix/V\). These are gross occupied-band rates. Pilot overhead is deducted only in the stationary net-rate calculation. Valid endpoints therefore imply the instance-wise bound \((V(h)-\fix(h))/(S\ln 2)\leq G\).

\subsection*{Theorem 1: a full-domain matching certificate}
We combine continuous-power bounds for individual pairs across a matching. Supplementary Note~2 gives the proofs, including vertex-price repair and signed rounding.

A power price \(\lambda\geq0\) lets us evaluate each pair separately while accounting for the total power budget. Assume \(0<p_-\leq p_+\) and \(P\geq Kp_-\). For a nonempty channel set \(\mathcal H\), let \(\mathcal Q=\mathcal H\times[p_-,p_+]^2\) contain the candidate channels and both users' feasible powers. Use \(H_{e,s}\) for the upper score of pair \(e\) on tone \(s\), and \(D(H)\) for the largest sum of table entries over feasible labelled schedules. Given tone prices \(\beta_s\in\mathbb R\), set \(w_e=\max_s(H_{e,s}-\beta_s)\) and let \(M(w)\) denote the maximum weight of a perfect user matching. For scores satisfying the first line below, the second line bounds every \(h\in\mathcal H\):
\begin{equation}
\begin{aligned}
 H_{e,s} &\geq \sup_{(h',a,b)\in\mathcal Q}
       \bigl\{q_{e,s}(a,b;h')-\lambda(a+b)\bigr\},\\[2pt]
 V(h)&\leq\lambda P+D(H)\leq\lambda P+\sum_{s=1}^{S}\beta_s+M(w).
\end{aligned}
\label{eq:upper}
\end{equation}
For the first rate inequality in equation~(\ref{eq:upper}), summing the pair bounds restores a power charge of at most \(\lambda P\). Each feasible schedule uses every tone once, so subtracting and restoring the tone prices preserves its score. Maximizing each adjusted entry over tones gives the second inequality.

The checker verifies this last maximization using free vertex potentials \(\alpha_i\) and nonnegative prices \(z_B\) on a family \(\mathcal B\) of odd vertex sets~\cite{edmonds}. Write \(\mathcal B_e=\{B\in\mathcal B:e\subseteq B\}\). For every edge \(e=\{i,j\}\), the checker requires
\[
 \alpha_i+\alpha_j+\sum_{B\in\mathcal B_e}z_B\ \geq\ w_e .
\]
A perfect matching uses each vertex once and contains at most \(\lfloor |B|/2\rfloor\) edges inside \(B\). Summing the inequalities thus bounds its weight by \(\sum_i\alpha_i+\sum_{B\in\mathcal B}\lfloor |B|/2\rfloor z_B\). Equality certifies that the proposed matching maximizes the priced edge weights \(w_e\). The corresponding labelled schedule can still fall short of the physical optimum. When this witness is unavailable, any vertex prices give a bound after setting \(\beta'_s=\max_{e=\{i,j\}}(H_{e,s}-\alpha_i-\alpha_j)\). Summing over a feasible schedule then gives \(U=\lambda P+\sum_i\alpha_i+\sum_s\beta'_s\). Vertex and tone prices may have either sign.

\subsection*{Theorem 2: label conflict and four-term accounting}
Let \(m_\beta\) attain \(M(w)\), \(R_\beta=\sum_s\beta_s+M(w)\), and \(A_\beta\) be the highest score of \(m_\beta\) when its pairs must occupy distinct tones, computed by bipartite assignment~\cite{kuhn}. Let \(d(\beta)=\max_e\{\max_s(H_{e,s}-\beta_s)-\min_s(H_{e,s}-\beta_s)\}\). Since \(A_\beta\) is feasible and \(R_\beta\) is an upper bound,
\begin{equation}
\begin{aligned}
0 &\le R_\beta-D(H)\le\dc=R_\beta-A_\beta,\\
\dc &=\sum_{e\in m_\beta}\bigl(w_e-H_{e,s(e)}+\beta_{s(e)}\bigr)
       \le S\,d(\beta).
\end{aligned}
\label{eq:conflict}
\end{equation}
The last inequality follows edge by edge. If \(H_{es}=a_e+b_s\), choosing \(\beta_s=b_s\) makes every edge range zero. Thus distinct labels cost nothing on a separable table, although the power enclosure can still contribute slack. In the four-user example, the two orthogonal disjoint matchings imply lower bounds \(4\ln2\pm(\beta_2-\beta_1)\) on \(R_\beta\). No edge can earn more than \(2\ln2\), so equal prices attain their common lower bound. Exhaustive evaluation of the six feasible assignments gives the stated positive discrepancy. More generally, minimizing the tone-price ceiling allows mixtures of matchings whose tone counts balance on average. Supplementary Note~2 derives this characterization and bounds the change in the discrepancy when channel scores are perturbed.

To decompose the interval for the delivered action, take the \emph{implemented} matching \(m\), its bijection \(s(e)\) and lower pair rates \(l_e\), all evaluated on the selected compatible table. For \(e=\{i,j\}\), write \(p_e=p_i+p_j\) and \(L=\sum_{e\in m}l_e\). Each \(l_e\) bounds the rate of the stored filters at the implemented powers on the channel event under consideration. For this decomposition the upper witness must satisfy \(U\geq\lambda P+\sum_s\beta_s+\sum_{e\in m}w_e\). The four terms are
\begin{equation}
\begin{aligned}
 T_1 &= U-\lambda P-\sum_s\beta_s-\sum_{e\in m}w_e,\\
 T_2 &= \sum_{e\in m}\bigl(w_e-H_{e,s(e)}+\beta_{s(e)}\bigr),\\
 T_3 &= \lambda\Bigl(P-\sum_i p_i\Bigr),\\
 T_4 &= \sum_{e\in m}\bigl(H_{e,s(e)}+\lambda p_e-l_e\bigr).
\end{aligned}
\label{eq:decomposition}
\end{equation}
The upper witness, the edge maxima, the power budget and the enclosure at the implemented powers each make the corresponding term nonnegative. The intermediate prices and scores cancel, leaving \(U-L=\sum_{j=1}^4T_j\) exactly. The assignment defining \(A_\beta\) belongs to the price-selected matching \(m_\beta\). The delivered schedule can use another matching from the candidate pool, in which case its \(T_2\) measures the label penalty for that action. The fourth term compares a pairwise upper score with the implemented floor. It can include slack from the power and channel enclosures, the chosen pair powers and filters, and downward rounding. Global loss from choosing other partners is not identifiable from this term alone.

\subsection*{Theorems 3 and 4: achieved rates under calibrated error}
A fixed rule maps the observed pilots to a feasible schedule, powers and nonzero filters. Both floors use this same decision. Write \(d_i\) for a filter normalized to unit length and \(s=s(i),j=j(i)\) for its tone and partner. Given a channel centre \(\hat h\) and directional radii \(b_{ik}\), define
\[
 \ell_i(b)=\frac{p_i\pos{|d_i^H\hat h_{is}|-b_{ii}}^2}
                 {1+p_j(|d_i^H\hat h_{js}|+b_{ij})^2}.
\]
On the event \(|d_i^H(h_{ks}-\hat h_{ks})|\leq b_{ik}\) for the desired and interfering signal of every user, the sum of \(\ln(1+\ell_i)\) rounded down is \(L_P\leq\fix(h)\). The reverse triangle inequality bounds the desired amplitude from below, and the ordinary triangle inequality bounds the interfering amplitude from above. On the full-vector event \(\|h_{ks}-\hat h_{ks}\|\leq\epsilon_{ks}\), Cauchy--Schwarz gives the same formula with \(b_{ik}=\epsilon_{ks}\), yielding \(L_F\), and Theorem 1 gives \(U_F\). Hence, on their intersection (Theorem 3),
\[
       \max(L_F,L_P)\ \leq\ \fix(h)\ \leq\ V(h)\ \leq\ U_F .
\]
Whenever the queried directions leave an unobserved subspace, a channel component in that subspace can increase the oracle rate without changing those projections. We therefore use the full-vector channel set when bounding the oracle.

For Theorem 4, condition on the fit and policy-selection data. The calibration blocks and the next test block must be exchangeable under this conditioning. The complete pilot-to-action map, including repairs and tie-breaking, and the positive scale functions \(\sigma_{ks},s_{ik}\) are fixed before calibration on data separate from the calibration blocks. Each block supplies two scores,
\[
\begin{aligned}
 T_F &= \max_{k,s}\frac{\|h_{ks}-\hat h_{ks}\|}{\sigma_{ks}},\\
 T_P &= \max_i\ \max_{k\in\mathcal K_i}
         \frac{|d_i^H(h_{k,s(i)}-\hat h_{k,s(i)})|}{s_{ik}},
 \qquad \mathcal K_i=\{i,j(i)\}.
\end{aligned}
\]
The scale functions may depend on a block's observed pilots through this frozen map. For \(a\in\{F,P\}\), we use the \(\lceil(n+1)(1-\alpha_a)\rceil\)-th smallest calibration score, appending infinity at rank \(n+1\). Multiplying the two thresholds by their corresponding scales defines the vector and projection radii. For exchangeable calibration and test blocks, random tie-breaking makes the rank of the test score uniform. Including all ties can only increase coverage~\cite{lei}. Each event has marginal probability at least \(1-\alpha_a\), and the union bound gives joint coverage of at least \(1-\alpha_F-\alpha_P\) under arbitrary dependence among users, tones and scores. At \(n=256\), \(\alpha_F=\alpha_P=0.025\) gives rank 251 for each, and at \(n=16\) it gives rank 17. The guarantee averages over calibration and target blocks. Coverage conditional on a particular calibration sample or on the scheduler accepting the interval can differ. An action selected after inspecting calibration residuals or the resulting radii requires separate calibration, or a block maximum covering every candidate action. If only a later noisy pilot is observed, transport to a true payload channel requires an additional error bound and its own error probability.

The Gaussian benchmark fixes all powers at one. The other simulations use the limits given above. Conditional on the specified channel centres, independent proper complex Gaussian errors have covariance \(\nu^2 I_{16}\). Set \(\nu^2=1/32\), \(A_0^2=160\), and let the two orthonormal directions have entries of magnitude \(1/4\). Each full score covers \(m=512\) independent links. Its squared normalized norm has gamma shape \(d=16\). Each projected score covers \(m=64\) independent scalar projections with \(d=1\). Their independence in this model follows from orthogonality and isotropic complex Gaussian errors. For \(F_d(x)=1-e^{-x}\sum_{j=0}^{d-1}x^j/j!\), the block cumulative distribution function (CDF) is \(F_d(q^2)^m\). At rank 251 among 256 blocks its CDF value has law \({\rm Beta}(251,6)\). For \(T\) with this law the physical radius has distribution \(\epsilon=\nu\sqrt{F_d^{-1}(T^{1/m})}\). For unit powers the per-tone upper and lower functions in bit/s/Hz are
\[
 u(x)=2\log_2(1+(A_0+x)^2),\qquad
 l(x)=2\log_2\!\left(1+\frac{\pos{A_0-x}^2}{1+x^2}\right).
\]
We integrated \(u(\epsilon_F)-l(\epsilon_F)\) and \(u(\epsilon_F)-l(\epsilon_P)\) over the radius distributions. Two quadratures in different coordinates agreed within \(2\times10^{-8}\). A 144-fractional-bit rational enclosure, using monotone bin bounds and an explicit Erlang tail, bounded the gaps in \([3.151125,3.156294]\) and \([1.333129,1.337516]\) bit/s/Hz. The quadratures give 3.1537 and 1.3353 bit/s/Hz to four decimal places. The wider intervals enclose the expectations rigorously.

\subsection*{Building and checking the interval}
The input is a complex \(K\times N\times S\) array with uncertainties in whitened units. We round channel centres to a common dyadic grid with integer components bounded by \(2^{22}\). For \(N\leq256\), each real or imaginary Gram sum fits the exact integer range through \(2^{53}\). We enlarge the supplied vector radius by an outward bound \(\sqrt{2N}\,\delta/2\) on the Euclidean rounding displacement, where \(\delta\) is the grid spacing. For centres \(a,b\) and radii \(u,v\), the magnitude of their inner product is at least \(\pos{|a^Hb|-\|a\|v-\|b\|u-uv}\). Dividing its square by the product of upper squared norms gives a lower bound on the squared normalized overlap when that product is positive. A zero product gives overlap zero. Bounding (\ref{eq:sinr}) with upper desired power and gain, lower interfering power and gain, and lower correlation encloses every realization in a power rectangle. We subtract the minimum power charge, and additionally take the better of that enclosure and a verified bound using tangents to the two individual concave log terms with a separately bounded correlation penalty.

Before refining any entry, we bound the full two-user power rectangle for every pair and tone. A change in price can reuse an unpriced enclosure. For each edge the two highest candidate tones enter a refinement shortlist. No more than \(\max(2,\lfloor K^2/32\rfloor)\) entries are fully divided into 16 rectangles. After all children cover a complete partition, the maximum child upper bound may tighten the parent. The stored bound is the minimum of that maximum and the parent bound. All other parents are retained. The priority score affects which entries tighten within the available time. The binary64 implementation uses nearest-even rounding with outward next-representable adjustments, gradual underflow and disabled fast-math reassociation. It range-reduces logarithms by powers of two, uses outward tabulation at \(1+j/256\) and a positive atanh series with a bounded remainder. Out-of-range or nonfinite data are rejected. The outward bounds certify the upper endpoint. Supplementary Note~7 and Supplementary Table~15 specify the numerical kernel and distinguish optional-leaf counts from total enclosure evaluations.

Each table entry is rounded \emph{up} to an integer multiple of \(2^{-40}\) nats before forming weights. The excess from rounding a schedule of \(S\) edges is less than \(S/2^{40}\) nats. A C++14 implementation using LEMON 1.3.1~\cite{lemon} exports a perfect matching, vertex prices and positive odd-set prices. An independent integer checker verifies the matching, every inequality and exact objective equality in the solver's quarter-tick potential units. The distinct-tone assignment and its potentials are also checked. A failed matching witness invokes the vertex-price repair described above. The stored powers are in units \(1/(20\cdot2^{20})\) W. Integer projection enforces the individual power limits and total budget. The complex filter \(v_i\) is retained unnormalized. Directed dot products enclose both signal terms and the noise contribution \(\|v_i\|^2\) in its SINR denominator. Integer decomposition numerators, their common denominator, per-pair floors, selected leaves, actions, prices and input/table hashes identify precisely what was checked.

\subsection*{Small-instance continuous and discrete references}
We evaluated the labelled-table optimum by integer dynamic programming on the same rounded table used by the certificate. We checked smaller tables by exhaustive enumeration. For the continuous-power reference, we enumerated all labelled schedules on 512 analytic channel blocks under the same power limits and total budget at \(K=2,4,6,8\). An outward power-box search kept the best feasible lower bound and the largest remaining upper bound, with a budget of 192 subdivisions per instance and a target width of \(10^{-6}\) bit/s/Hz. The median remaining widths were \(5.8534\times10^{-5}\), \(7.3372\times10^{-5}\) and \(1.5043\times10^{-4}\) bit/s/Hz at \(K=4,6,8\), respectively. These medians are approximately three orders of magnitude below the 0.15-bit/s/Hz scheduling criterion. The reference brackets nevertheless have long tails, reaching 0.514, 0.854 and 0.929 bit/s/Hz, respectively. All 512 intervals at each size exceed the finer target, so we report both endpoints.

For two users, a separate calculation resolves the continuous problem. With total power 0.20 W and positive channel gains, scaling both powers upward increases both signal-to-interference-plus-noise ratios until the sum constraint is tight. The feasible boundary is therefore \(p_2=0.20-p_1\), \(0.05\leq p_1\leq0.15\) W. Exact rational Gram coefficients reduce stationarity to a polynomial equation. We isolated its real roots in intervals of width at most \(2^{-75}\), included both boundary endpoints and evaluated outward logarithmic bounds at all candidates. We converted each interval width from nats before rounding the displayed rate endpoints. The largest two-user width was \(5.13\times10^{-15}\) bit/s/Hz. We also recomputed the stored-filter rate from exact rational inner products on every small analytic and O1 instance (Supplementary Note~7). This separate root calculation supplies the two-user continuous reference, whereas the integer dynamic program supplies the labelled optimum for a specified upper-bound table.

\subsection*{Price search, feasible schedules and learned proposals}
We set the deterministic power price from the median single-user marginal rate. We evaluate seven tone-price proposals through 32 users and six at 64 users, starting with zero tone prices. Subsequent anchors include centred mean scores and the tone potentials of a relaxation with two user slots per tone. To update the most promising price vector, we count excess favourite-tone requests from matched pairs and from individual user neighbourhoods. We weight the centred assignment potentials by 0.58 and the previous prices by 0.42. The demand correction uses weights of 0.75 for pairs and 0.25 for neighbourhoods, with a decreasing step starting at 0.05. These settings control how quickly the search tightens the interval. The same checker accepts each resulting bound.

We collect the sequential pairing, matchings from the successive prices and up to six promising disjoint edges per tone. Distinct-tone assignments, exchanges of tone labels, and exchanges of partners between two pairs generate feasible alternatives. For up to three promising schedules, we attempt four projected gradient updates. Each gradient is normalized by its largest absolute component, with a \(10^{-12}\) floor, and multiplied by \(0.05\) W and successive relative steps \(1,1/2,1/4,1/8\). Each proposal is projected onto the stored feasible powers and accepted if the nominal pair-rate sum does not decrease. The filters are then recomputed for those powers. We retain the largest verified floor and the smallest verified ceiling. We evaluate the four terms using that ceiling's table and that floor's action.

To propose starting prices, we use a network of width 32 with three residual message-passing layers. Each pair--tone factor receives its current state, the pair and tone embeddings, and the sum and absolute difference of its two user embeddings. Means across tones and pairs update the corresponding nodes. Incident pair messages update each user. Heads propose pair--tone scores, powers, centred prices, a power price, box anchors and priorities. Training combines balanced positive/negative schedule loss, squared errors for powers, prices and anchors, and relative-priority losses. The tone potentials of the optimal assignment for the teacher's selected schedule are centred to give the price target.

Three seeds (1729, 2718, 3141) each saw 4,096 blocks for 15 epochs with Adam (step \(10^{-3}\), size-homogeneous batches up to eight, gradient norm cap 3). Fit data comprised clustered-delay-line (CDL) profiles A, B and C and analytic channels at 8, 16 or 32 users. Four-site urban-macrocell (UMa) and 64-user examples were withheld from fitting. The mean validation fraction below gap thresholds 0.02, 0.05, 0.10 and 0.15 on 256 other blocks selected seed 1729. The learned warm start adds three price proposals to the full deterministic search. In the exploratory control, we replace these proposals with three further deterministic updates. All methods use the same channels, with rotating order over three calls. The rerun reproduced all 1,024 deterministic gaps. Neural inference and candidate evaluation take additional time, so we also report complete runtimes. A separate comparator uses one forward pass. Both learned routes retain full-domain verification. Supplementary Note~4 gives the layer equations, losses and solver pseudocode. In the combined route, learning adds price proposals to the deterministic table. Learned power prices, tangent anchors and leaf priorities affect the separate standalone comparator. On a Quadro RTX 4000 Max-Q, the synchronized standalone forward pass had 95th-percentile latencies of 15.65 ms at 32 users and 65.63 ms at 64 users.

\subsection*{Channels, uncertainty and measurement processing}
The native experiments used Sionna 2.2.0, with \texttt{spec\_version='19.2'} for CDL and UMa, selecting the 3GPP TR 38.901 version 19.2.0 parameter tables~\cite{tr,sionna}. At 3.5 GHz the three single-site profiles used a 16-element array and 100-ns delay spread. UMa used four 4-element arrays, at height 25 m on the corners of a 500-m square, with outdoor users at 1.5 m sampled uniformly within the central 420-m square. Path loss and shadowing were enabled. Sixty-four bins at 312.5 kHz spacing spanned 20 MHz. Nested user counts and evenly spaced labelled bins came from the same realization. A common block scale set the median per-antenna signal-to-noise ratio (SNR) at 0.1 W to the target value before cropping. All comparisons therefore use normalized channels. The analytic controls used nearly parallel users, frequency-flat and frequency-selective channels, and a benchmark with orthogonal partners on every assigned tone. The two ray-traced scenes used base station 1, the O1 60-GHz and ASU 3.5-GHz source files, and 64 disjoint guarded tiles each, avoiding pilot tiles and their immediate neighbours. We considered \(-5,0,10,20\) dB and 16, 32 or 64 users in the wider sensitivity grid (13,824 cells). The two primary endpoints used the native 10-dB channels at 32 and 64 users. Relative to the operating sum rates of each regime, these absolute gap criteria correspond to roughly 1\% of the certified ceiling on known channels (0.15 bit/s/Hz), 2.6--3.3\% of the median ceiling under noisy pilots at 10 dB (0.50 bit/s/Hz), and 0.41--0.42\% of the median early lower rate on the six-user hardware captures (0.10 bit/s/Hz).

To interpret differences between the two ray-traced scenes, we examined the stored user vectors and the base-station ray records for the same held-out users. For each tile we averaged \(r\) over distinct user pairs and tones. On each tone we also formed the Gram matrix of unit-norm user vectors and computed its effective rank \(\exp(-\sum_j t_j\ln t_j)\), where \(t_j\) are its eigenvalues divided by their sum and \(0\ln0=0\). We averaged ranks over tones before summarizing tiles. As a separate measure of sensitivity to partner choice, we held both powers at 0.10 W and used equation~(\ref{eq:sinr}) to evaluate each user's rate with every possible partner. We averaged the difference between the 90th and 10th percentiles over users and tones within each tile. The two scene comparisons use medians across tiles. The ray files supply path counts, received path powers, base-station departure directions and a direct-path indicator for each link.

In the imperfect-channel experiments, independent complex Gaussian pilot noise was added to the normalized channels with variance per complex component \(1/E_p\), where \(E_p=32\) or \(32\,10^{\mathrm{SNR}_{\rm dB}/10}\). The same frozen pilot-only least-squares scheduling policy was used for calibration and testing. For each family, SNR and pilot-energy convention we used 256 independent whole blocks for calibration. Each family contributed a further 128 test topologies. Utility required \(L_P>0\), \((U_F-L_P)/(S\ln2)\leq0.50\), \(L_P/U_F\geq0.95\), and elapsed time at most 1 s. This interval included pilot-to-action processing, bound construction and the simulation residual-score calculations performed within the timed routine. We checked joint coverage against the known simulated channels.

The stationary cost calculation compares the zero-price, equal-power-initialized matching control, including its feasible power refinement, with learned warm start on the same least-squares estimate. It uses 32 routine pilot symbols per 256-symbol block, horizon \(N_{\rm blk}\in\{1,4,16,64,256,1024\}\), and \(M\in\{0,16,64,256,1024\}\) additional 32-symbol refresh blocks for the warm action. Four error events across two actions share a 5\% budget (rank 254 of 256). For each fixed action, projection-radius triangle inequalities bound its desired and interfering amplitudes in both directions, producing rate intervals \([l_c,u_c]\) and \([l_w,u_w]\) in bit/s/Hz. The payload fractions are \(\eta_c=1-32/256\), \(\eta_w=\max(0,\eta_c-32M/(256N_{\rm blk}))\). A feasible budget has \(32M+32N_{\rm blk}\leq256N_{\rm blk}\). We classify a gain only if \(\eta_w l_w-\eta_c u_c>10^{-5}\) bit/s/Hz. A loss requires \(\eta_wu_w-\eta_cl_c<-10^{-5}\) bit/s/Hz. Other affordable cases are indeterminate. These net gains assume stationary per-block rates. Coverage for a sequence of future blocks would require joint calibration over that horizon.

The distributed Rice recordings used a 3.6 GHz carrier and a 64-point fast Fourier transform (FFT) over 5 MHz. We took signed data bins \(-13,+1,+14\) (positions 12, 24 and 36 in zero-based increasing signed order across 48 data tones). The first 128 of 2,001 frames were reserved for quality, timing and noise estimation. Two nominal windows \([176,240)\) and \([256,320)\), plus the stored per-user alignment corrections, supplied disjoint pilots in each user slot. Channels were divided by the corresponding noise standard deviation per antenna and tone without simulated block scaling. The radio-frequency (RF) stream mask required at least one user with at least 6 dB signal-to-noise and adjacent-pilot coherence 0.60. All six decision pilots had to pass detection.

The distributed HUB line-of-sight and colocated HUB non-line-of-sight catalogue entries have identical filenames, byte lengths and acquisition timestamps for both chains. All 24 common sampled byte ranges for chain A and all 22 for chain B had matching hashes. We therefore treated the two directory entries as probable copies and attributed the records to their catalogue labels. In an additional nested-size replay, the distributed LO chain-B entry had no qualifying calibration pilots, so the scheduled windows were recorded as abstentions. All six users transmitted simultaneously in the recorded payload. The early and later pilots permit an offline evaluation of the proposed two-user schedule.

For the additional chronological RENEW analysis, we fitted the error radii on the first 256 eligible frames after the original 128-frame preprocessing prefix. The pilot-to-action map remained fixed. Each calibration score was either the largest vector norm of the LTS2-minus-LTS1 difference over users and tones or the largest absolute projection of those differences along the unit-normalized deployed filters, including both desired and interfering users. Each event received probability budget 0.025, giving order-statistic rank $\lceil0.975(256+1)\rceil=251$. We froze these two radii and evaluated the remaining 1,614 HUB and 1,513 LO frames. The separate CPU checker replayed every calibrated certificate. Our target in this replay is the observed later pilot. We can evaluate the geometric bounds on these frames, but temporal dependence leaves the exchangeability assumption for population coverage unverified. On every eligible early input we also enumerated all 90 labelled pairings at fixed equal powers, using the optimal linear-combining rate of each pair. Exhaustive discrete enumeration does not optimize continuous powers.

For the phase-invariant replay, the projected score is the maximum of \(\big||d_i^Hh_k^{(1)}|-|d_i^Hh_k^{(0)}|\big|\) over the 12 scheduled desired/interference terms, and the vector score is \(\max_{k,s}\min_\theta\|e^{-{\rm i}\theta}h_{ks}^{(1)}-h_{ks}^{(0)}\|_2\). If the magnitude residual is bounded by \(b_{ik}\), the desired magnitude is at least \((|d_i^Hh_i^{(0)}|-b_{ii})_+\), and the interfering magnitude is at most \(|d_i^Hh_j^{(0)}|+b_{ij}\). Substitution into the fixed-combiner SINR therefore gives the projected rate-floor formula above; phase-aligned vector balls also bound the oracle because every candidate pair rate is invariant to independent user/tone phases.


For Argos, the first 64 frames fixed the RF mask. We selected 32 later, approximately evenly spaced frames per recording and used FFT windows starting at sample offsets 47 and 111. All eight channels were retained. Each of six users had to supply finite, nonzero vectors on 52 occupied tones, occupied-to-null concentration of at least 3 dB and adjacent-pilot coherence at least 0.50. We used a common mask per recording and fixed all frame choices before measuring user correlation. Automatic gain control was enabled in the source, so the statistic refers to recorded digital vectors.

\subsection*{Detection of the recorded six-user transmission}
We decoded the distributed chain-A data slots separately from the paired-schedule pilot replay. All six transmitters occupied slot 47. We used all 64 stored receive-stream positions. The first 512 physical frames of each capture fixed one amplitude gain per user. Frames 512--2,000 were held out. The file repeats one transmitted data frame, so both calibration and evaluation span the same six fixed symbol sequences. We used only known pilot tones to estimate phase and gain.

Each slot had ten 80-sample OFDM symbols with a 16-sample cyclic prefix. Our 64-sample FFT windows started at offsets \(172+80j\), \(j=0,\ldots,9\). The first window of each user's training slot supplied its column of the channel matrix after division by the known long training sequence. We used the Moore--Penrose inverse on each of 52 occupied tones and held the spatial weights fixed over the data slot. For each user and OFDM symbol we formed \(c=\tfrac14\sum_{k\in\mathcal P}z_k p_k^*\), where \(z_k\) is the equalized value on control tone \(k\) and \(p_k\) is its known transmitted pilot. We multiplied the 48 equalized data values by \(\exp[-\mathrm{i}\arg(c)]\) and divided by that user's median \(|c|\) across eligible calibration frames and all ten OFDM symbols.

We required every user's early training-window mean squared magnitude across samples and streams to reach twice the larger of one squared ADC count and its mean over the first 64 guard samples. A zero data slot or a smallest channel singular value at or below \(10^{-6}\) of the largest on any occupied tone also caused abstention. We checked the centred FFT order on all six time-domain transmitter files and matched both sign bits of all 480 frequency-domain data symbols per user to their reference labels. The 48 data tones excluded the four control pilots. Each hard decision assigned the signs of the real and imaginary parts to the high and low label bits. Symbol and bit error rates divided the error counts by scored symbols and uncoded bits. RMS error vector magnitude was \(100\sqrt{\sum|\hat x-x|^2/\sum|x|^2}\) percent. A noise-free six-stream superposition and a one-symbol rotation of the scoring reference served as controls. We report descriptive counts within each acquisition.

\subsection*{Statistical analyses and reproducibility}
For the retrospective seven-ceiling analysis, four primary methods each contributed the smallest verified ceiling over three runs. Three controls each contributed one ceiling. For each topology we summed the median call duration of each primary method and the single duration of each control. We also summed all 15 call durations. Neither sum measures the elapsed time of an integrated implementation.

For the two primary endpoints we ran each algorithm three times per topology and size on one workstation, processing one array at a time and rotating the algorithm order. We measured wall time from entry to the scheduling routine, after channel cropping and normalization, through serialization of the certificate. We report the 95th percentile across all calls. A topology passed only when the largest certified gap over its three calls was at most 0.15 bit/s/Hz and the median wall time met the 250-ms deadline at 32 users or the 1,000-ms deadline at 64 users. At 32 and 64 users the required joint pass fractions were 0.80 and 0.60, and the fixed one-sided 97.5\% lower limits had to exceed 0.70 and 0.50. Each bootstrap replicate resampled 128 \emph{whole topologies} with replacement inside each native profile, preserving all pairing across methods (10,000 draws, seed 19073). The design fixes the four profile counts. The pooled Clopper--Pearson limits use a common binomial success probability as a reference. We report percentile bootstrap limits, using the 2.5th percentile for the one-sided 97.5\% lower limits and the 2.5th and 97.5th percentiles for two-sided 95\% intervals. After both primary endpoints passed, we tested the secondary learning hypothesis. It specified an absolute gain of at least ten percentage points at 32 users, a positive lower confidence limit and no reversal in the sign of the gain at 64 users. We also report the exact two-sided McNemar test, conditioning on the number of discordant topology pairs and using a binomial probability of one half for each direction under the null hypothesis. Missing, rejected or late decisions count against the relevant utility denominator.

We computed Spearman rank correlations and resampled whole topologies within each native profile for their two-sided 95\% percentile intervals and specified a lower confidence limit above 0.30 for the posterior metric. For the nine-level selectivity series, each block used a fixed sequential pairing and equal powers. A block crossed when the price-selected matching's posterior conflict exceeded half the certified interval for that fixed action. A sustained transition required at least half the blocks to cross at two adjacent levels. Resampling preserved each block's entire series. The acquisition is the unit of replication for comparing clock configurations. Conditional on each fitted calibration threshold, we used a lower-tailed exact binomial test at level 0.025 to screen the independent simulation test blocks for a joint hit probability below 0.95. This diagnostic conditions on the fitted thresholds. The exchangeability result averages over the calibration sample as well.

We fixed the configuration before evaluating the held-out test cohort. The protocol had no external registration. The source/checkpoint manifest is timestamped 13 November 2026 at 14:38:50 UTC and the linked pilot lock at 14:44:37 UTC; the project repository records both timestamps, their hashes and the protocol in \nolinkurl{code/manuscript/revision/protocol_provenance.json}. The original primary thresholds were fixed before held-out evaluation. The manifest records when the revisions were made. The subsequent iteration-matched, time-budget, depth-sweep and phase-invariant controls are exploratory. Their timing-only pilot uses 16 hash-selected topologies, three calls and rotating order to choose 38 and 63 deterministic price iterations; these budgets are then fixed for the full 512-topology comparison at each size. Pass criteria and three-repeat aggregation remain unchanged, and every returned certificate is replayed. We matched median time across methods. Individual topologies could take different times.


The computation used an Intel Core i9-10885H at 2.40 GHz (8 physical cores, 16 logical processors), 32 GB installed RAM and 64-bit Windows 11 (build 26200). The CPU environment used Python 3.13.5, NumPy 2.3.4, SciPy 1.16.3, Numba 0.62.1, NetworkX 3.4.2, Zig 0.13.0 and LEMON 1.3.1. BLAS and OpenMP were limited to one thread for the verification runs. The GPU was an NVIDIA Quadro RTX 4000 with Max-Q Design. The four methods, two sizes and 512 native topologies produced 12,288 timed calls. Repeated identical results reduce to 8,192 distinct serialized certificates. Three seeds required 0.293 GPU training hours in total. The confirmatory runs and analysis took 0.86 workstation hours, excluding the later CPU verification and manuscript preparation. Timing sensitivity on 52 topologies selected by the frozen hash rule gave median shuffled/original time ratios of 1.70 and 1.55 for deterministic search at 32 and 64 users. Simulation jobs shared the workstation. The repository documents input reconstruction and proof replay.

We used a separate CPU implementation, independent of training, to authenticate all 8,192 certificates and their 12,288 timing references and to verify every four-term identity. Sampling without replacement using the release hash as a seed selected 58 of 1,152 eligible units. All 232 exact 16-user integer objectives across four SNRs agreed, and exhaustive enumeration agreed on 120 nested smaller tables. The audit reproduced the primary outcomes and recorded no numerical violations. Supplementary Table~16 summarizes the independent checks and intentional corruption controls. The checker ran on the same workstation using the recorded inputs and numerical conventions above.

